\section{Recovering the polarized ramification variety}
\label{sec:polarization}
Let $A=\OO_Y(1)$ and $B=\OO_{\Pj(S)}(1)$. In the graph description
\eqref{eq:graph-E} write
\[
 Z=Y\times\Pj(S),\qquad W=\mathcal H^*\otimes Q,
 \qquad Q=(S\otimes\OO)/\Lambda,
 \qquad \pi:E=\Tot_ZW\to Z,
\]
and let $\rho:E\to Y$ be the composite projection. Constant
one-dimensional vector-space factors in determinant lines below are
suppressed; they change none of the isomorphism classes or section rings.

\begin{lemma}[No extra fibrewise functions]\label{lem:no-functions}
The natural map $\OO_Y\to\rho_*\OO_E$ is an isomorphism.
\end{lemma}
\begin{proof}
The assertion is local on $Y$, where $\mathcal H$ can be trivialized.
If $H_0$ is its fibre space, the affine bundle projection gives
\[
 \pi_*\OO_E=\bigoplus_{d\ge0}\Sym^d(H_0\otimes Q^*).
\]
We claim that
\begin{equation}\label{eq:vanish-symmetric}
 H^0(\Pj(S),\Sym^d(H_0\otimes Q^*))=0\qquad(d>0).
\end{equation}
The injection $Q^*\hookrightarrow S^*\otimes\OO$ embeds a global
section in the constant space $\Sym^d(H_0\otimes S^*)$.
Such an element is a homogeneous polynomial on $\Hom(H_0,S)$.
At every line $\Lambda\subset S$, membership in
$\Sym^d(H_0\otimes(S/\Lambda)^*)$ says this polynomial factors
through $\Hom(H_0,S/\Lambda)$; it is invariant under translations
by $\Hom(H_0,\Lambda)$. Taking the lines of a basis of $S$ shows
invariance under every translation of $\Hom(H_0,S)$. In characteristic
zero such a polynomial is constant, so its positive-degree part is zero.
This proves \eqref{eq:vanish-symmetric}. The degree-zero part is $\C$.

The same calculation tensors with the coordinate ring of each affine
trivializing open of $Y$: cohomology on the projective space commutes
with this flat extension from $\C$, and with the displayed direct sum
in degree zero. Hence the pushforward is precisely $\OO_Y$, with
its natural multiplication. The local identifications glue.
\end{proof}

\begin{proposition}[The canonical nilradical line]\label{prop:canonical-line}
As line bundles on $E$,
\begin{equation}\label{eq:two-lines}
 \NN\simeq\pi^*(A^{-1}\boxtimes B^{-1}),\qquad
 \omega_E\simeq\pi^*(A^{-(p-1)}\boxtimes B^{-(p+e-1)}).
\end{equation}
Consequently
\begin{equation}\label{eq:line-cancellation}
 \omega_E\otimes\NN^{\otimes-(p+e-1)}\simeq\rho^*A^e.
\end{equation}
\end{proposition}
\begin{proof}
The determinant defining $\widehat\Delta$ is a section of
$(\det\mathcal K)^*\otimes\det V$. On $E$ the exact sequence
$0\to\Lambda\to\mathcal K\to\mathcal H\to0$ gives
$\II_\Delta|_E=\det\mathcal H\otimes\Lambda$, up to a constant
factor. Since $\det\mathcal H=A^{-1}$ and $\Lambda=B^{-1}$,
Theorem~\ref{thm:normal-form} gives the first formula.

The determinant of $W$ is
$(\det\mathcal H^*)^{p-1}\otimes(\det Q)^{e-1}
=A^{p-1}\boxtimes B^{e-1}$.
The canonical bundle of a vector-bundle total space is
$\pi^*(\omega_Z\otimes\det W^{-1})$, as follows by taking the
determinant of its relative cotangent sequence. Here
$\omega_Z=\OO_Y\boxtimes B^{-p}$ because $\omega_Y=\OO_Y$.
This proves the second formula. Subtracting the two exponents gives
\eqref{eq:line-cancellation}: the $B$-exponent vanishes and the
$A$-exponent is $-(p-1)+(p+e-1)=e$.
\end{proof}

\begin{proof}[Proof of Theorem~\ref{thm:main-reconstruction}]
An isomorphism of failure schemes preserves the nilradical and its
annihilator, hence $E$, its invertible module $\NN$, and its canonical
bundle. It therefore preserves $\LL=\omega_E\otimes
\NN^{-(p+e-1)}$ and its graded section ring.
Proposition~\ref{prop:canonical-line}, Lemma~\ref{lem:no-functions},
and the projection formula give
\[
 H^0(E,\LL^n)=H^0(Y,A^{en}\otimes\rho_*\OO_E)=H^0(Y,A^{en})
\]
compatibly with multiplication. The full section ring of the ample
line bundle $A^e$ recovers $(Y,A^e)$ by $\Proj$. In this case it is
also the $e$th Veronese of the homogeneous coordinate ring of a smooth
hypersurface, so there is no ambiguity about the grading line bundle.

For a K3 surface the Picard group has no torsion. Here is the short
argument needed: a nontrivial torsion line bundle $T$ and its inverse
have no nonzero sections, since a nonzero effective torsion divisor
has zero intersection with an ample divisor, which is impossible.
Riemann--Roch and Serre duality would then give
$2=\chi(T)=-h^1(T)\le0$, a contradiction. Thus an isomorphism
carrying $A^4$ to $(A')^4$ carries $A$ to $A'$. This recovers the
quartic polarization, and its complete linear system recovers the
quartic embedding up to projective equivalence.

Finally Proposition~\ref{prop:full-fitting} recovers $D$ from
$\widehat D$ by restricting to the smooth locus of the reduction.
This operation is intrinsic and is preserved by every abstract scheme
isomorphism. The preceding proof applies to that restriction.
\end{proof}

\begin{remark}
For $e=3$ the theorem recovers $(Y,A^3)$ directly. A chosen cubic root
in the Picard group need not be unique. As polarized genus-one curves,
however, the possible degree-three roots are isomorphic: after choosing
an origin, the map $y\mapsto t_y^*A\otimes A^{-1}$ onto
$\Pic^0(Y)$ is multiplication by three, hence surjective. No choice
of an origin is needed for the stated section-ring reconstruction.
For $e=4$ the torsion-free argument gives uniqueness without translations.
\end{remark}

\begin{remark}[Relation to the earlier birational argument]
The compactification $\Pj_Z(\OO\oplus W)$ is birational to
$Y\times\Pj^{e(p-1)}$. Thus the earlier cancellation proof is a
special case of the maximal rationally connected quotient: for a
non-uniruled base, that quotient of a product with projective space
is the base, birationally. See the rational-quotient construction in
\cite[Chapter IV, \S5]{Kollar}. This explains the earlier abstract
recovery, which is retained in the companion. The present proof does
not rely on that argument: the section ring gives the actual projective
model and its line bundle, not merely a birational class.
\end{remark}
